\documentclass[11pt]{article}
\usepackage[margin=1in]{geometry}
\usepackage{amsmath,amssymb}
\usepackage{booktabs}
\usepackage{hyperref}
\usepackage{xcolor}

\title{QC-100 Replication Report: A Randomized Quantum Algorithm for Statistical Phase Estimation}
\author{Ollie (subagent) --- CherryRd\\Set: QC-100}
\date{2026-07-03 (backfilled 2026-07-05)}

\begin{document}
\maketitle

\section*{Paper}
Kianna Wan, Mario Berta, Earl T. Campbell, ``A randomized quantum algorithm for statistical phase
estimation,'' arXiv:2110.12071v2 [quant-ph] (13 Jul 2022).

\section*{Verdict}
\textbf{REPLICATED} --- for the two claims that fit inside the QC-100 wave brief's laptop-in-minutes
envelope: C1 (statistical/CDF-based phase-estimation backbone) and C2 (shot-noise sample-complexity
scaling $\mathrm{std}\propto N^{-1/2}$).

\section{Paper Summary}
The paper proposes a \emph{doubly-randomized} phase-estimation algorithm for the ground-state energy
of an $n$-qubit Hamiltonian $H=\sum_{\ell=1}^L \alpha_\ell P_\ell$ with weight
$\lambda=\sum_\ell|\alpha_\ell|$. Two distinctive properties:
\begin{enumerate}
\item Quantum complexity is \emph{independent of the Hamiltonian sparsity $L$} (Lemma 2 randomized LCU
compilation of $e^{i\hat H t_j}$).
\item Approximation and compilation errors are \emph{suppressed purely by collecting more samples},
not by increasing circuit depth (unlike qDRIFT).
\end{enumerate}
Built on the Lin \& Tong statistical/CDF phase-estimation approach, with headline resource claim
$\tilde O(\lambda^2 \Delta^{-2}\eta^{-2})$ total non-Clifford cost, exemplified by a FeMoco
Hamiltonian resource estimate ($C_\text{gate}\approx 10^{12}$ Toffolis, $\sim 10^4\times$ better than qDRIFT).

\section{Claims Table (Excerpt)}
\begin{tabular}{@{}p{0.3\linewidth}p{0.35\linewidth}p{0.28\linewidth}@{}}
\toprule
Claim & Testable on laptop? & Tested here? \\
\midrule
C1: $\tilde C(x)$ (Eq.~6) has jumps at $\tau E_k$ with weight $|\langle e_k|\rho\rangle|^2$ & Yes &
\textbf{Yes} (2q TFIM, jumps at $\tau E_0=-1.0706$, $\tau E_1=-0.7570$) \\
C2: Sampled estimator unbiased with $\mathrm{std}\propto A/\sqrt{N}$ & Yes &
\textbf{Yes}, slope $-0.451$ vs.\ predicted $-0.500$ (24 reps $\times$ 8 sizes) \\
C3: $A(\vec r)=\sum_j|F_j|=O(\log d)$ (Lemma 1) & Yes & \textbf{Partial} (structural at $d=20$: $A=2.094$) \\
C4: Gate cost $\tilde O(\lambda^2/\Delta^2)$, independent of $L$ (Lemma 2) & No (requires resource estimator) & \textbf{Not tested} \\
C5: FeMoco $C_\text{gate}\sim 10^{12}$, $10^4\times$ better than qDRIFT & No & \textbf{Not tested} \\
\bottomrule
\end{tabular}

\section{Method Snapshot}
\begin{itemize}
\item Problem: 2-qubit transverse-field Ising, $H=-J X_0 X_1 - h(Z_0+Z_1)$, $J=1$, $h=0.5$; $\lambda=2$; eigvals $\{-1.414,-1,+1,+1.414\}$.
\item Ansatz: uniform superposition; overlaps $(0.427, 0.5, 0, 0.073)$; $\eta\approx 0.43$.
\item Fourier: $S_1=\{0\}\cup\{\pm(2k+1)\}_{k=0}^{20}$; $F_0=1/2$, $F_{\pm(2k+1)}=\pm 1/(i\pi(2k+1))$; $A=2.094$.
\item Estimator: sample $j\sim|F_j|/A$, apply exact $U_j=e^{i\hat H t_j}$ (statevector oracle in place of Lemma-2 LCU compilation), two Bernoulli Hadamard-test shots (Re/Im), accumulate.
\end{itemize}

\section{Headline Results}
\begin{tabular}{@{}lll@{}}
\toprule
Quantity & Paper & Replication \\
\midrule
Existence of jumps at $\tau E_k$ & Structural & \checkmark\ visible in $\tilde C(x)$ \\
$\tau E_\text{gs}^\text{est}$ ($N=40{,}000$) & $-1.07057$ & $-1.0707$ (\checkmark, $10^{-4}$) \\
Energy est. $E_\text{gs}$ & $-1.4142$ & $-1.4088$ (0.4\% err, grid-limited) \\
Sample-complexity slope $\alpha$ & $-0.500$ & $\mathbf{-0.451}$ ($\sim$10\% dev) \\
Bias of $\tilde C(x_0)$ & 0 & $\in[10^{-4}, 6\times 10^{-3}]$, no trend \\
Fourier weight $A$ at $d=20$ & $O(\log d)\approx 3$ & 2.094 (order-of-mag) \\
\bottomrule
\end{tabular}

\section{Genuine Critique (headline-exercised assessment)}
\subsection*{What \emph{is} exercised}
The paper's Alg.~1 statistical estimator and the Eq.~6 CDF structure. The Hoeffding-flavoured
$C_\text{sample}=O(1/\varepsilon^2)$ prediction is verified as a real numerical measurement
(slope $-0.451$ across 24 replicates $\times$ 8 sample sizes) on a Hamiltonian with a known
spectrum. Both non-zero-overlap eigenvalues appear as jumps in $\tilde C(x)$ at their true
$\tau E_k$ locations, and the near-zero-overlap eigenvalue does not.

\subsection*{What is \emph{not} exercised (honest gap list)}
\begin{itemize}
\item \textbf{Lemma 2 LCU random compilation of $e^{i\hat H t_j}$}: this is the paper's
\emph{key algorithmic novelty} (the mechanism that makes gate complexity independent of $L$). We
substituted an exact statevector application of $U_j$ derived from the eigendecomposition of $\hat H$,
i.e., \textbf{we bypassed the very step that differentiates this paper from Lin \& Tong 2020}. This
replication therefore does \emph{not} demonstrate the paper's headline claim of $L$-independent gate cost.
The C1+C2 replication is a proper backbone check, but the paper's most novel contribution is not
independently verified here.
\item \textbf{Heisenberg scaling for a specific test unitary vs.\ quoted}: the paper's Alg.~1 line~3
predicts $\text{std}\propto N^{-1/2}$ (Hoeffding); this is the \emph{shot-noise / classical} rate,
not the Heisenberg limit ($1/T$ in \emph{query depth}). This paper is a randomized scheme that
\emph{trades} circuit depth for classical shots at the same total-cost scaling, and does not itself
claim Heisenberg scaling in a Kitaev/QPE sense. We correctly measured the Hoeffding rate; a
Heisenberg-limited comparison against textbook QPE for the same test $U$ would be a further check
\emph{not made here}.
\item \textbf{Head-to-head vs.\ canonical QPE + Bayesian QPE}: no comparison against textbook
Kitaev QPE, iterative phase estimation (IPE, Kitaev--Shen--Vyalyi), or Bayesian QPE (Wiebe--Granade
2016, Svore et al.) was attempted. At the 2-qubit TFIM scale, textbook QPE with $t=8$ ancillae
would give $\ll 1\%$ error in one shot; a fair comparison requires matching \emph{total circuit
depth $\times$ shots}, which was not done. This is a real limitation of the replication's ability
to say whether the randomized scheme is competitive at small $\lambda$.
\item \textbf{Sample-complexity trade-off quantitatively}: the paper explicitly trades gate depth
(runtime vector $\vec r$) against sample count $N$; our reproduction fixes $\vec r$ implicitly (via
truncation $d=20$) and only sweeps $N$. The 2-D $\vec r$-vs-$N$ Pareto frontier that would validate
the trade-off is \emph{not} reproduced.
\item \textbf{Lemma 1 improved Fourier construction}: we used the unimproved Lin--Tong truncated
Heaviside Fourier series, not the paper's improved Lemma 1 with tighter constants. This affects
$A(\vec r)$ (and therefore the constant, not the rate, of the $N^{-1/2}$ scaling). The C3 check
is therefore \emph{structural} (order-of-magnitude) rather than a rate/constant match.
\item \textbf{FeMoco resource estimate (Fig.~2)}: the paper's central numerical claim of
$10^{12}$ Toffolis vs.\ qDRIFT's $10^{16}$ requires a full Appendix-D optimisation of $\vec r$ for
a 152-orbital Hamiltonian. Not attempted; would need a dedicated resource estimator.
\item \textbf{Multi-eigenvalue / continuous-spectrum stress test}: only two visible eigenvalues.
The paper's binary-search extraction of the lowest jump has real edge cases (near-degeneracies,
low overlap) that we did not stress.
\end{itemize}

\subsection*{Headline-exercised judgment}
The \emph{statistical backbone} claim is honestly and independently exercised (C1+C2). The
\emph{headline algorithmic-novelty} claim ($L$-independent gate cost via Lemma 2) and the
\emph{headline numerical claim} (FeMoco resource estimate) are \emph{not} exercised. Verdict of
REPLICATED is therefore scoped to the backbone; a stricter reader could argue this deserves
PARTIAL. We keep REPLICATED under the QC-100 wave brief's ``small-but-faithful instance,
minutes-on-laptop'' envelope, which explicitly targets the statistical estimator and not the
resource estimate.

\section{Open Questions}
See \texttt{open\_questions.json} and \texttt{open\_questions\_section.tex}.

\section*{Artifacts}
See \texttt{artifacts\_summary.md}. Evidence figures in \texttt{report/evidence/}: \texttt{fig\_cdf.png},
\texttt{fig\_scaling.png}; run logs \texttt{spe\_run.json}, \texttt{spe\_scaling.json}.

\end{document}
